% v20: current support consequences. Five necessary types; no existence claim.
\subsection{残る五型と二価頂点の高近傍}
\label{sf20:section}

\begin{corollary}[非空支持の連結性と二価頂点の制約]
\label{sf20:no-leaves}
正の局所整数 charts，marked puncture の整数 clutching，対向 $D$ の相殺，
および node 非負性を満たす，五次の標準 $K_5$ branch graph を考える．
高分母 support $H=\{ij:q_{ij}>1\}$ が非空なら，
$H$ は五頂点全てを使う二点連結 graph であり，最小次数は二以上である．
さらに
\[
 |E(H)|\ge7,\qquad b_1(H)\ge3
\]
であり，各二価頂点の二つの高近傍は，ともに次数四を持つ．
特に二価頂点は高三角形に属する．
高辺が七本なら $H$ は book 型 $G_{5,13}$ である．
二価頂点の個数は零，一，三のいずれかであり，
一つの場合は $G_{5,15}$，三つの場合は book 型である．
一つの core の高分母 ray が高々一本なら，全ての分母は一である．
\end{corollary}

\begin{proof}
系~\ref{kg17:low-core-propagation} により，非空の場合に孤立頂点はない．
木成分，単純閉路成分，二価二頂点を持つ三角形，葉付き四角形，
$K_{2,3}$，diamond，二葉三角形および四使用頂点 $K_4$ の障害を適用すると，
五頂点上の十一型が残る．
定理~\ref{hs18:no-house}，\ref{gten18:no-support}，\ref{lh18:no-K4-leaf} により，
$G_{5,5},G_{5,10},G_{5,12}$ が除かれる．
この段階で葉を持つ型，または六辺の型は $G_{5,7},G_{5,9}$ だけであり，
定理~\ref{g719:no-G5-7} と定理~\ref{dl19:no-G5-9} がそれらを排除する．
従って最小次数二以上，辺数七以上，$b_1(H)=|E(H)|-4\ge3$ を得る．
定理~\ref{g51120:no-support} により $G_{5,11}$ も除かれる．
次の五代表の表から，二価頂点の両近傍が次数四を持つことと，
七辺の場合および二価頂点数についての主張が従う．

二点連結性は最小次数と辺数からも直接従う．切断頂点 $v$ があるなら，
$H-v$ の各成分は最小次数二の条件により二頂点以上を持つ．
残る四頂点は二頂点ずつの二成分に分かれ，各成分の両頂点は互いに接続し，
ともに $v$ に接続しなければならない．
これは二つの三角形が $v$ だけを共有する六辺の graph であり，辺数七以上に反する．
最後の結論は，孤立頂点と葉がいずれも許されないことの言い換えである．
\end{proof}

\begin{corollary}[五型六十六標識への帰着]
\label{sf20:remaining-support}
前系の仮定の下で，非空の高分母 support の必要候補は次の五型に限られる．
標識を保つ候補は六十六個である．$\overline H=K_5\setminus H$ とし，
補 graph の表示では孤立頂点を省略する．
\begin{center}\small
\begin{tabular}{crcrcr}
型&$|E|$&降順の次数列&$b_1$&$\overline H$&標識付き個数\\\hline
$G_{5,13}$&7&$(4,4,2,2,2)$&3&$K_3$&10\\
$G_{5,14}$&8&$(4,3,3,3,3)$&4&$2K_2$&15\\
$G_{5,15}$&8&$(4,4,3,3,2)$&4&$P_3$&30\\
$G_{5,16}$&9&$(4,4,4,3,3)$&5&$K_2$&10\\
$G_{5,17}$&10&$(4,4,4,4,4)$&6&$\varnothing$&1
\end{tabular}
\end{center}
識別子は固定し，欠番を詰めない．
これらは必要条件を通過した graph の個数であり，
整数接着データ，作用 atlas，または Lagrangian fibrations の存在数を表さない．
\end{corollary}

\begin{proof}
全 $2^{10}=1024$ 個の標識付き部分 graph に既述の障害を適用する．
$G_{5,11}$ は自己同型群の位数が二であるため，その排除は六十標識を除く．
五代表の自己同型群の位数は表の順に $12,8,4,12,120$ であり，
軌道数 $5!/|\operatorname{Aut}(H)|$ の和は
\[
 10+15+30+10+1=66
\]
となる．各代表の補 graph を取ると第五列が得られる．
\nolinkurl{verify_support_agenda_v20.py} は全 graph に対して，
固定した識別子と軌道数，各頂点の削除後の連結性，
二価頂点の二近傍の次数を確認する．二価頂点の確認は全標識を通して六十件となる．
結果は \nolinkurl{support_agenda_verification_v20.json} に全標識とともに保存する．
この有限 graph の列挙は，先行する整数接着の排除定理を適用するものであり，
それらの証明の代わりではない．
\end{proof}

残る低辺は高々三本であり，三本なら低三角形を作る book 型だけである．
book の低三角形には定理~\ref{bk19:finite-completion} の固定 core に対する有限完成法がある．
さらに定理~\ref{bk20:no-all-negative} と系~\ref{bk20:mixed-necessary} により，
正の枝と負の枝を混在させる必要がある．
定理~\ref{bd20:nonzero-cycle} は，全ての低 $D$ が非零で，
負の端から正の端へ向けた三本の低辺が有向三角形をなすことを要求する．
これらは book 全体の不存在をまだ与えない．

$G_{5,15},G_{5,16},G_{5,17}$ は $K_4$ を部分 graph として含む．
四使用頂点 $K_4$ や $K_4+\mathrm{leaf}$ の排除は support 全体に関する定理なので，
これらの密な graph を同じ理由で除くことはできない．
三高分母の同時非負 $D$ 問題は，定理~\ref{hd20:nonsingular} によって
非特異な Schur 領域へ帰着されるが，その領域での一般問題は未解決である．
五候補の整数接着・幾何的実現と，一般の五次分母剛性も未解決である．
